\section{可执行检查究竟确认了什么}
\label{sec:audit}
现在将这些区别用于第~\ref{subsec:execution-sample}~节的 14 对选定候选，以及第~\ref{subsec:constructed-materials}~节另行介绍的构造材料。编译与公理检查确认形式文件的一些性质；定义反例和布尔对照则进一步检验，这些性质是否足以保证完成原任务。

\subsection{编译、目标提取、公理与题意}
\label{subsec:lean-checks}
一个 Lean 文件可以包含多个有名称的定义和定理。\emph{指定目标}是评价前选定的声明。同一文件中的另一条定理证明成功，不能替代对这个目标的检查。

需要区分四个问题。第一，文件能否在记录的软件环境中\emph{通过编译}？这会检查语法、名称、类型及所提供的证明在该环境中是否被接受。第二，能否找到并提取指定目标？第三，这个声明依赖哪些\emph{公理}？公理是作为推理起点接受的命题；Lean 不仅能报告直接依赖，还能追踪经由其他定义和定理形成的依赖。第四，这个被接受的声明是否表达并满足任务原本要求的数学内容？前三项为第四项提供证据，却不能单独决定答案。特别是，内核检查的是\emph{实际写出的形式命题}是否具有相对于其定义和假设的有效证明；这个命题与任务原意是否一致，还需要另行解释和核对\cite{leanaxioms,leanvalidation}。

历史候选执行采用 Lean 4.31.0，每个候选最多运行 240 秒，同时最多运行两个进程。实验使用两套公理政策。第一套排除 \code{sorryAx}，它是与未完成证明的占位写法相关的依赖；这种写法允许文件保留一个尚未补齐的证明。第二套只允许 \code{propext}、\code{Classical.choice} 和 \code{Quot.sound} 这三个列明的基础公理。通过某套政策，只说明满足了相应的公理依赖规则，并不说明定理的假设符合任务要求，也不说明定义采用了任务指定的约定。这两套政策是实验检查，不能当作完整历史审查要求的新旧版本。

28 个历史候选版本全部尝试了执行：17 个编译成功，11 个失败，没有超时。17 个成功文件中，16 个能够提取预先指定的目标，并且这 16 个目标全部通过两套公理政策。另一个成功文件 \code{ex\_3\_1\_2/from} 缺少指定声明，核查没有在看到结果后改选一个方便通过的目标。因此，只有六对候选拥有完整的四格公理评价表，且各表四格均为一。这些结果描述的是本次记录的当前环境。原历史环境的完整依赖没有重建，因此当前编译失败不能证明文件在原历史环境中也曾失败。

\subsection{针对直接假设结论的狭窄检查}
\label{subsec:assumed-conclusion}
一个形式上有效的证明，可能证明了比原任务更弱的命题。用初等例子说明：假设任务要求证明，对每个实数 $x$ 都有 $x^2\geq 0$。如果改写成“若 $x^2\geq 0$，则 $x^2\geq 0$”，证明就很容易，因为假设已经给出了结论。这个条件命题在逻辑上成立，却没有证明原任务要求的无条件结论。这只是帮助理解的说明例，不是一条新增的历史观察。

另一条要求检查，正是询问主定理是否把它的\emph{完整结论}直接列为公开假设。此项检查由单个人工智能评估者完成。它不能发现所有弱化命题的方式，也不能发现所有间接把证明负担移入假设的方式；本研究没有独立专家给出的参考判定。14 对候选为本项要求评估提供 $14\times4=56$ 个评价单元，其中27 个可测量，23 个不适用，四个不确定，两个无效；只有六张表完整且可比较。“不适用”表示这个特定检查不适合该对象；“不确定”表示证据不能给出确定判断；“无效”表示拟议评价不能按所声明的方式使用。它们都没有被悄悄换算成一次测得的失败。这些数量说明哪里可以完成测量，不说明历史数学内容有多大比例正确。

\subsection{四种评价方法与诊断可答性}
\label{subsec:diagnostics}
对一对候选和两套要求，每张表提出六个诊断问题：分别在两套要求下，更换候选有什么变化；分别对两个候选，更换要求有什么变化；候选与要求同时更换时总共有什么变化；以及候选变化的表现是否随所用要求而改变。最后一个问题，就是前面数学框架介绍的交互项。

四种方法的区别在于允许使用的信息。\emph{历史标签}使用保存的审查判断；若不能把原评价器与本次评价器视为同一个，就暂不给出答案。\emph{固定 $R_1$}用第二套要求检查两个候选，其中 $R_1$ 是第二套要求的记号。\emph{四格}使用两个候选与两套要求的全部四种组合。\emph{四格加夏普利分配}使用同样的观察，再按照前文介绍的办法，对两种更换顺序取平均，从而分配对角总变化。

固定要求的方法同样可以使用已核实的有效规则相同信息，以及成立的逻辑蕴含。例如，通过更严格的公理允许清单必然意味着通过更宽松的清单，那么一次严格政策的通过结果，就足以推出宽松政策也通过，无须再次执行。这种格子被记录为逻辑推导。如果刻意不向简单方法提供这条信息，就会不公平地夸大四格方法的信息优势。

\begin{table}[htbp]
\centering\small
\caption{六个诊断问题中，能够给出有依据的单一数值答案的数量。“可答位置”是同一行所有方法共同使用的分母。}
\label{tab:availability}
\begin{tabularx}{\linewidth}{Y r r r r r}
\toprule
证据组 & 可答位置 & 历史标签 & 固定 $R_1$ & 四格 & 加夏普利 \\
\midrule
构造例子（7 个） & 42 & 4 & 17 & 26 & 26 \\
公理：开发组（6 对） & 36 & 0 & 13 & 13 & 13 \\
公理：留出组（8 对） & 48 & 0 & 27 & 27 & 27 \\
要求：开发组（6 对） & 36 & 0 & 12 & 12 & 12 \\
要求：留出组（8 对） & 48 & 0 & 25 & 25 & 25 \\
\bottomrule
\end{tabularx}
\end{table}

例如，第一行包含七张构造表，每张提出六个问题，因此有 $7\times6=42$ 个可能的答案位置。四格方法能够回答其中的 26 个问题；这不是对 42 个独立抽样任务作出了 26 次正确判断。开发组的每行使用六对历史候选，所以分母为 $6\times6=36$；留出组的每行使用八对，因此分母为 $8\times6=48$。公理检查和要求检查重复使用同一批历史候选对，不能把这些位置当成独立样本。历史标签列中的零，表示另一评价器保存的标签不能提供这里规定的当前测量，并不表示所有历史标签都是错的。

在构造例子中，四格方法比固定 $R_1$ 多回答九个问题；在各历史组中没有增加可答问题。夏普利分配没有增加可回答的问题，因为它没有增加新的观察。构造表另有 26 个可核对、预先手写的参考答案，没有发现实现结果与之不一致。这些参考答案用于检查设计案例的算术，不是经独立裁决的历史语义正确性基准。历史组中差值全为零应怎样作统计解释，将在本报告的统计部分另行讨论。

\subsection{为什么空集能区分两个定义}
\label{subsec:definition-witness}
考虑历史任务 \code{def\_2\_1}。两个集合之间的\emph{双射}，就是把双方元素一一配对，每个元素都恰好有一个配对对象，既不遗漏，也不重复。自然数集 $\mathbb N=\{0,1,2,\ldots\}$ 有无限多个元素。有限集合不能与全部自然数一一配对。特别地，空集 $\varnothing$ 没有任何元素，连自然数 $0$ 所需要的配对对象都无法提供。

“可数”有两种常见约定。一种把有限集合以及能与自然数集建立双射的集合都包括在内；另一种专指能与自然数集建立双射的集合。保留的两版任务文字都使用第二种约定。较早的候选却采用 Lean 中的 \code{A.Countable}，把有限集合也包括进来。较新的候选使用了任务要求的双射约定，并另外命名了范围更广的“至多可数”概念。因此，问题在于是否忠实遵守\emph{本任务的约定}，并不是说较宽的约定通常就是错误的。

核查所用的\emph{参考谓词}是：“存在从 $A$ 到整个自然数集的双射。”这里的 $A$ 是输入给这个定义的集合；谓词就是对一个对象提出的真假断言。为便于查阅实现，这条参考断言在 Lean 中写成
\[
\code{Nonempty}\bigl(A\simeq(\code{Set.univ}:\code{Set}\ \mathbb N)\bigr).
\]
符号 $\simeq$ 表示要求一个双射；\code{Set.univ} 表示指定类型的全部元素组成的集合，这里就是整个自然数集；\code{Nonempty} 断言至少存在一个这样的双射。它\emph{不是}仅仅断言 $A$ 中有元素。把任务的自然语言翻译成这条形式参考，仍然是本次核查中由人工智能作出、可以公开检查的解释，没有独立专家参考；不能只凭编译通过就认为这一步解释已经得到认证。

空集现在就给出了完整的反例论证。旧定义接纳空集，因为其约定包括有限集合。参考定义拒绝空集，因为任何双射都必须让 $0$ 在空集中找到一个原像，而空集中没有元素。新定义也因同样的理由拒绝空集。只需一个反例，就足以推翻“两个谓词对\emph{所有}集合都给出相同答案”的断言。

见证文件把两版候选的主体保留在不同命名空间，共用原有的 Mathlib 导入，然后补入六条已证明的命题：旧谓词接纳空集；参考谓词拒绝空集；新谓词拒绝空集；新谓词对所有集合都与参考谓词一致；旧谓词并非对所有集合都与参考一致；以及新旧谓词确实不同。新谓词与参考的一致性可通过展开定义得到。这六条命题共同描述一个被调查的案例，不能当作六次独立的语义核查成功。

成功编译的退出码为零，用时 217.685 秒。两版候选定义打印出的公理集合都为空；四条见证命题只依赖 \code{Quot.sound}，另两条的公理集合为空。来源与消息验证器通过了 23 项检查。此前一次因 Windows 路径过长引起的启动失败发生在编译之前，仍保留在档案中；成功运行使用了较短目录。这些证据确认的是相对于所选参考的这一项数学差异，不是对整个修改后文件或全部下游定理的认证。

\subsection{对一个布尔任务穷尽行为检查}
\label{subsec:boolean-audit}
有限域挑战不要求读者理解高等数学。任务有两个布尔输入，每个输入只能取假或真；当且仅当\emph{恰好一个}输入为真时，函数必须返回真。输入组合总共只有 $2\times2=4$ 种，因此表~\ref{tab:truth-reference} 可以完整列出所有要求。参考表在执行之前就已固定。

\begin{table}[htbp]
\centering\small
\caption{双输入任务的完整参考真值表。}
\label{tab:truth-reference}
\begin{tabularx}{\linewidth}{c c c Y}
\toprule
第一个输入 & 第二个输入 & 要求输出 & 理由 \\
\midrule
假 & 假 & 假 & 没有输入为真。 \\
假 & 真 & 真 & 只有第二个输入为真。 \\
真 & 假 & 真 & 只有第一个输入为真。 \\
真 & 真 & 假 & 两个输入都为真，不是恰好一个为真。 \\
\bottomrule
\end{tabularx}
\end{table}

实验有意设计了九种实现：三种正确写法、五种具有相同函数类型的错误写法，以及一种编译失败的对照。函数类型只规定“接受两个布尔值，返回一个布尔值”，说明允许输入和输出哪类值，却不说明某个具体输入应该返回\emph{哪个}值。因此，“永远返回假”的函数可以与正确的异或函数拥有相同类型。

\begin{table}[htbp]
\centering\small
\caption{有限任务保存的观察结果。八种可执行实现均通过公理检查和打印类型检查。横线表示没有可用的执行证据。}
\label{tab:finite}
\begin{tabularx}{\linewidth}{l Y c c}
\toprule
编号 & 实现 & 不匹配输入数 & 行为结果 \\
\midrule
01 & 异或 & 0/4 & 一致 \\
02 & 布尔值不相等比较 & 0/4 & 一致 \\
03 & 用条件分支实现异或 & 0/4 & 一致 \\
04 & 合取：两个输入都为真 & 3/4 & 不一致 \\
05 & 析取：至少一个输入为真 & 1/4 & 不一致 \\
06 & 经辅助函数实现相等比较 & 4/4 & 不一致 \\
07 & 恒定返回假 & 2/4 & 不一致 \\
08 & 只返回第一个输入 & 2/4 & 不一致 \\
09 & 调用未定义的函数 & -- & 不确定 \\
\bottomrule
\end{tabularx}
\end{table}

表~\ref{tab:finite} 将每种可执行实现与全部四行参考逐一比较。例如，析取只在两个输入都为真时出错：它返回真，而任务要求假。因此，穷尽执行区分出了三种一致实现和五种不一致实现；无法执行的版本仍记为不确定。通过前两种检查不足以确认行为符合任务。由于九种实现是围绕\emph{同一个任务}有意设计的，“八个中有五个”不能解释为某个总体的误接纳率，也不能视为八个独立抽样任务的证据。按具体能力组织行为测试已有方法上的先例\cite{checklist}，但这里没有照搬该研究在自然语言处理中的性能结论。有限输入上的执行检查，也不同于上一个定义例中由内核检查的数学命题证明。

原评价器将每种实现与一个固定的正确基线分别编译，共有 18 次候选编译；另有九次行为探针编译，合计 27 次尝试。重复编译的基线属于测试安排，不会增加独立证据。一个单独的验证器从保存的标准输出重建观察，不调用原来的评分计算。它在修正换行规范化缺陷后通过了 166 项检查：最初九个测试文件的哈希失败来自 LF 与 CRLF 换行差异。两版验证器和首次失败报告均已保留。规范化仅用于测试文件的换行，实际执行源码和探针的哈希仍要求完全一致。没有通过重新做实验来替换那次失败的验证记录。

\subsection{结果重建与矛盾输入的拒绝}
\label{subsec:reconstruction}
排除运行耗时之后，保存的 35 张比较表均被精确重建。早期重建测试集通过 52 项测试；其来源与数量核对验证器，对读取的 227 份来源快照通过了 6,093 项检查。这些结果检验的是既定实现规则下的可复现性，不是独立的语义正确性。

核查还发现了一个输入一致性问题。假设输入声称，新候选在两列面对\emph{完全相同的有效要求}，但实测格子却是 $q_{10}=0$、$q_{11}=1$。这里，第一格是新候选在第一套要求下的结果，第二格是它在第二套要求下的结果。对同一候选执行相同的确定性检查，不应同时得到拒绝与接纳。旧基线可能根据“要求相同”的声明推得要求变化为零，而直接相减却得到 $1-0=1$。便携实现现在会在写出结果之前拒绝这种矛盾输入，不再把两个互相冲突的答案都呈现为有效结果。

该实现的 50 项单元与回归测试，和先前的 52 项重建测试是两套不同测试。本报告保留这一实现及其 35 个兼容性案例。本轮写作修订改变的是说明方式，没有改变已有实验、实验结果或生产形式化语料。
